\chapter{Fundierung und allgemeine Satzschemata}

\section{Ausschlüsse durch Fundierung}

\begin{theorem}[Kein Element gehört sich selbst an]
\label{b03guide:thm:noself}
\[
 \forall A\,(A\notin A).
\]
Insbesondere gilt für die Russell-Teilmenge aus
Theorem~\ref{b03guide:thm:russell} nun zusätzlich $R_A=A$.
\end{theorem}

\begin{theorem}[Keine gegenseitige Mitgliedschaft]
\[
 A\in B\Rightarrow B\notin A.
\]
Außerdem können $A\in B$, $B\in C$ und $C\in A$ nicht
gleichzeitig gelten.
\end{theorem}

\begin{theorem}[Ein Glied ohne gemeinsame Elemente]
Jede nichtleere Menge $A$ enthält eine Menge $a$ mit
\[
 a\cap A=\varnothing.
\]
Insbesondere gibt es keine nichtleere Menge $A$, für die
$a\cap A\ne\varnothing$ für alle $a\in A$ gilt.
\end{theorem}

\begin{theorem}[Mitgliedschaft, Inklusion und Vereinigung]
\[
\begin{gathered}
 A\subseteq B\Rightarrow B\notin A,\qquad
 A\in B\Rightarrow B\nsubseteq A,\\
 A\notin\bigcup A.
\end{gathered}
\]
\end{theorem}

\begin{theorem}[Hinzufügen der Menge als neues Element]
Für jede Menge $A$ gelten
\[
 A\subsetneq A\cup\{A\},\qquad
 (A\cup\{A\})\setminus A=\{A\}.
\]
Für beliebige Mengen $A,B$ gilt
\[
 A\cup\{A\}=B\cup\{B\}\Longleftrightarrow A=B.
\]
Die Mengenausdrücke werden hier ohne einen Nachfolgerbegriff
für Zahlen verwendet.
\end{theorem}

\begin{theorem}[Potenzmenge und Mitgliedschaft]
\[
 \mathcal P(A)\notin A,\qquad \mathcal P(A)\ne A.
\]
\end{theorem}

\section{Mengeninduktion}

\begin{theorem}[Schema der Mengeninduktion]
\label{b03guide:thm:setinduction}
Für jede Formel $\varphi(x)$ der Mengensprache mit festen
Mengenparametern gilt
\[
 \Bigl[\forall x\,\bigl(
       (\forall y\in x\,\varphi(y))\Rightarrow\varphi(x)\bigr)\Bigr]
      \Rightarrow\forall x\,\varphi(x).
\]
Der Satz betrifft unmittelbar die Elemente jeder Menge.
\cite{miquel}
\end{theorem}

\begin{theorem}[Fundierung für definierbare Eigenschaften]
Für jede Formel $\varphi(x)$ mit festen Mengenparametern gilt
\[
 \exists x\,\varphi(x)
 \Rightarrow
 \exists a\,\bigl(\varphi(a)\land
                    \forall y\in a\,\neg\varphi(y)\bigr).
\]
Die durch $\varphi$ beschriebenen Mengen müssen selbst keine
Menge bilden.
\end{theorem}

\section{Sammlung ohne Eindeutigkeitsforderung}

\begin{theorem}[Sammlungsschema]
\label{b03guide:thm:collection}
Für jede Formel $\varphi(x,y)$ mit festen Mengenparametern gilt
\[
 \forall x\in A\,\exists y\,\varphi(x,y)
 \quad\Rightarrow\quad
 \exists B\,\forall x\in A\,\exists y\in B\,\varphi(x,y).
\]
Die Menge $B$ enthält zu jedem $x\in A$ mindestens eine passende
Menge $y$. Zusätzlich kann verlangt werden:
\[
 \forall y\in B\,\exists x\in A\,\varphi(x,y).
\]
Für $A=\varnothing$ ist dann $B=\varnothing$.
Das Schema ist ein Satz von ZF ohne Auswahlaxiom; Eindeutigkeit
der Zeugen wird nicht gefordert. \cite{miquel,holland}
\end{theorem}

